\section{Background}
\label{sec:background}

\subsection{The LFP/graphite cell}
The cells are A123 APR18650M1A, a cylindrical 18650 cell with a lithium iron
phosphate (LiFePO$_4$, LFP) cathode and a graphite anode. The nominal
capacity is 1.1\,Ah and the voltage window is 2.0--3.6\,V. During
discharge, lithium ions leave the graphite and move into the LFP. During
charge they move the other way. LFP has a very flat voltage plateau (about
3.3\,V at low current). This makes the cell's state hard to infer from
voltage alone, but it also means that small changes in the shape of the
curve can carry information.

\subsection{C-rate and state of charge}
Current is expressed as a \emph{C-rate} relative to the cell's capacity. A
current of 1\,C fully charges or discharges the cell in one hour, which is
1.1\,A for this cell. So \Crate{4} means 4.4\,A, or a full charge in
15~minutes. \emph{State of charge} (SOC) is the stored charge as a
percentage of capacity. A higher C-rate charges faster but stresses the
cell more: it heats the cell, raises the overpotential, and makes lithium
plating on the graphite more likely.

\subsection{Degradation mechanisms}
Capacity fade in these cells is usually explained by three
modes~\cite{birkl2017}:
\begin{itemize}
  \item \textbf{Loss of lithium inventory (LLI):} lithium is used up in side
  reactions, mainly growth of the solid--electrolyte interphase (SEI) layer
  and lithium plating during fast charging.
  \item \textbf{Loss of active material (LAM):} particles crack, lose
  electrical contact or dissolve, so less electrode material can store
  lithium.
  \item \textbf{Increase in impedance:} the internal resistance rises,
  which increases the overpotential and reduces the capacity that can be
  used at high current.
\end{itemize}
Each mode changes the shape of the discharge curve $Q(V)$ in its own way.
This is why curve-based features can detect degradation early, before the
total capacity drops.

\subsection{Charge--discharge protocols}
A standard charging method is \emph{constant current--constant voltage}
(CC-CV). The cell is charged at a fixed current until it reaches the upper
voltage limit, then held at that voltage while the current decays. The
fast-charging policies studied here replace the first part of the charge
with \emph{two constant-current steps}, described in
Section~\ref{sec:policies}.

\subsection{End of life and cycle life}
\emph{Cycle life} is defined as the number of cycles until the discharge
capacity first falls below 80\% of nominal, i.e.\ 0.88\,Ah. This is the
quantity the models in this project predict.
